\section{Modul 08. Stack}

\subsection{Tujuan}

Setelah menyelesaikan modul ini, mahasiswa mampu:
\begin{enumerate}[nosep]
    \item menjelaskan sifat \emph{last in, first out} (LIFO) pada
          stack;
    \item menggunakan \texttt{PUSH} dan \texttt{POP} serta menjelaskan
          perubahan \texttt{SP};
    \item menjelaskan peran \texttt{BP} sebagai titik acuan parameter
          di dalam stack;
    \item menulis prosedur yang menerima parameter lewat stack;
    \item menjelaskan \texttt{RET} dekat tanpa angka dan pembuangan
          parameter oleh pemanggil;
    \item merakit, menjalankan, dan memeriksa stack dengan DEBUG.
\end{enumerate}

\subsection{Dasar singkat}

\subsubsection*{Stack dan sifat LIFO}

Stack adalah daerah memori untuk menyimpan data sementara. Puncak
stack ditunjuk pasangan \texttt{SS:SP}. Stack 8086 tumbuh ke alamat
yang lebih rendah. Satu sel stack berukuran satu word atau dua byte.

\texttt{PUSH} menyimpan satu word ke puncak stack: \texttt{SP}
dikurangi dua, lalu operand disalin ke \texttt{SS:SP}. \texttt{POP}
mengambil satu word dari puncak stack: isi \texttt{SS:SP} disalin ke
tujuan, lalu \texttt{SP} ditambah dua. Karena data terakhir yang
masuk berada di puncak, data itu pula yang pertama diambil. Sifat
inilah yang disebut LIFO. Urutan \texttt{POP} harus kebalikan urutan
\texttt{PUSH} bila nilai ingin kembali ke register asalnya.

\begin{table}[H]
    \centering
    \caption{Jejak \texttt{PUSH} dan \texttt{POP} pada program
    pertama. Nilai awal \texttt{AX=1111h}, \texttt{BX=2222h}.}
    \label{tab:m08-lifo}
    \small
    \begin{tabular}{@{}llll@{}}
        \toprule
        \textbf{Instruksi} & \textbf{SP} & \textbf{Puncak stack} &
        \textbf{Register sesudah} \\
        \midrule
        \texttt{PUSH AX} & berkurang 2 & 1111h & \texttt{AX=1111h} \\
        \texttt{PUSH BX} & berkurang 2 & 2222h & \texttt{BX=2222h} \\
        \texttt{POP CX} & bertambah 2 & 1111h & \texttt{CX=2222h} \\
        \texttt{POP DX} & bertambah 2 & kosong & \texttt{DX=1111h} \\
        \bottomrule
    \end{tabular}
\end{table}

\subsubsection*{Register SP dan BP}

\texttt{SP} (\emph{Stack Pointer}) menunjuk puncak stack yang aktif.
Nilainya berubah otomatis setiap kali \texttt{PUSH} atau \texttt{POP}
dijalankan. \texttt{BP} (\emph{Base Pointer}) dipakai sebagai titik
acuan data di dalam stack. Setelah prosedur menyimpan \texttt{BP} lama
dan menyalin \texttt{SP} ke \texttt{BP}, isi stack dapat dibaca
memakai bentuk \texttt{[BP+disp]} tanpa terpengaruh perubahan
\texttt{SP} berikutnya. Alamat \texttt{[BP+disp]} memakai segmen
\texttt{SS}.

\subsubsection*{Parameter prosedur lewat stack}

Pemanggil mendorong parameter ke stack sebelum \texttt{CALL}.
Parameter kedua didorong lebih dahulu, parameter pertama didorong
paling akhir. Prosedur lalu menyimpan \texttt{BP} lama, menyalin
\texttt{SP} ke \texttt{BP}, dan membaca parameter pada
\texttt{[BP+4]} serta \texttt{[BP+6]}. Susunan bingkai stack
ditunjukkan pada Tabel~\ref{tab:m08-frame}.

\begin{table}[H]
    \centering
    \caption{Susunan bingkai stack prosedur dengan dua parameter word.}
    \label{tab:m08-frame}
    \small
    \begin{tabular}{@{}ll@{}}
        \toprule
        \textbf{Alamat} & \textbf{Isi} \\
        \midrule
        \texttt{[BP+6]} & parameter kedua \\
        \texttt{[BP+4]} & parameter pertama \\
        \texttt{[BP+2]} & alamat kembali \\
        \texttt{[BP+0]} & nilai \texttt{BP} lama \\
        \bottomrule
    \end{tabular}
\end{table}

\texttt{RET} dekat tanpa angka hanya mengambil alamat kembali dari
stack. Parameter yang sudah didorong masih tertinggal di stack,
sehingga pemanggil membuangnya dengan \texttt{ADD SP,4} untuk dua
parameter word. Jumlah \texttt{PUSH} dan \texttt{POP} di dalam
prosedur harus seimbang agar \texttt{SP} kembali ke posisi semula.

Rujukan singkat: Gunawan, Bagian 12; S'to, Bab XIV; praktikum
\texttt{bab5.tex}; buku ajar, \texttt{chapters/14} bagian stack dan
\texttt{param-bp.asm}.

\subsection{Alat dan bahan}

\begin{itemize}[nosep]
    \item Komputer dengan emu8086 terpasang.
    \item DOSBox dengan TASM, TLINK, dan DEBUG.
    \item Berkas \texttt{code/m08-push-pop.asm} dan
          \texttt{code/m08-param-stack.asm}.
    \item Lembar jurnal praktikum.
\end{itemize}

\subsection{Contoh program}

Program pertama menyimpan \texttt{AX} lalu \texttt{BX} ke stack, dan
mengambilnya kembali ke \texttt{CX} lalu \texttt{DX}. Karena sifat
LIFO, \texttt{CX} memperoleh nilai \texttt{BX} dan \texttt{DX}
memperoleh nilai \texttt{AX}.

\lstinputlisting[
    caption={Menyimpan dan mengambil dua register dengan
    \texttt{PUSH} dan \texttt{POP}.},
    label={lst:m08-push-pop}
]{code/m08-push-pop.asm}

Program kedua memakai prosedur \texttt{TAMBAH} dengan dua parameter
word di stack. Hasil penjumlahan dikembalikan pada \texttt{AX},
kemudian dicetak sebagai satu digit.

\lstinputlisting[
    caption={Prosedur dengan parameter lewat stack dan
    \texttt{[BP+disp]}.},
    label={lst:m08-param}
]{code/m08-param-stack.asm}

\subsection{Langkah kerja}

\begin{enumerate}
    \item Siapkan folder kerja, lalu salin kedua berkas \texttt{.asm}
          ke dalamnya.
    \item Rakit dan tautkan program pertama. Jalankan kedua perintah
          berikut pada baris terpisah:
          \begin{quote}
          \texttt{tasm m08-push-pop.asm}\\
          \texttt{tlink m08-push-pop.obj}
          \end{quote}
    \item Jalankan \texttt{debug m08-push-pop.exe}, lalu ketik
          \texttt{r}. Catat nilai \texttt{SS} dan \texttt{SP}. Ketik
          \texttt{t} beberapa kali sampai kedua \texttt{PUSH}
          dijalankan, lalu periksa stack dengan
          \texttt{d SS:offset L 8} memakai offset \texttt{SP} yang
          tercatat. Lanjutkan \texttt{t} sampai kedua \texttt{POP},
          catat \texttt{CX}, \texttt{DX}, dan \texttt{SP}, lalu keluar
          dengan \texttt{q}.
    \item Rakit dan tautkan program kedua dengan perintah yang sama,
          lalu jalankan. Catat karakter yang tercetak.
    \item Jalankan program kedua dengan \emph{single step} pada
          emu8086. Catat \texttt{SP} dan \texttt{BP} di dalam
          prosedur, lalu periksa isi \texttt{[BP+4]} dan
          \texttt{[BP+6]}. Pastikan jumlahnya sama dengan keluaran.
    \item Ubah kedua parameter dari 3 dan 5 menjadi 2 dan 4. Rakit
          ulang, jalankan, dan pastikan keluarannya berubah.
\end{enumerate}

\subsection{Tugas}

\begin{enumerate}
    \item Jelaskan arti LIFO dan mengapa urutan \texttt{POP} harus
          kebalikan urutan \texttt{PUSH}.
    \item Jelaskan perubahan \texttt{SP} pada setiap \texttt{PUSH} dan
          \texttt{POP}.
    \item Jelaskan peran \texttt{BP} dan arti \texttt{[BP+4]} serta
          \texttt{[BP+6]} pada prosedur \texttt{TAMBAH}.
    \item Jelaskan mengapa pemanggil perlu menulis \texttt{ADD SP,4}
          sesudah \texttt{CALL} dan \texttt{RET}.
    \item Diketahui \texttt{AX=AAAAh} dan \texttt{BX=BBBBh}. Tentukan
          isi \texttt{AX} dan \texttt{BX} sesudah runtunan
          \texttt{PUSH AX}, \texttt{PUSH BX}, \texttt{POP AX},
          \texttt{POP BX}.
\end{enumerate}

\subsection{Hasil yang dikumpulkan}

\begin{enumerate}
    \item Kedua berkas \texttt{.asm} beserta versi modifikasinya.
    \item Tangkapan layar keluaran program kedua dan pemeriksaan stack
          dengan DEBUG.
    \item Tabel jejak berisi \texttt{SP}, isi puncak stack, dan isi
          register untuk kedua program.
    \item Lembar jurnal berisi jawaban langkah kerja dan tugas.
\end{enumerate}
